\documentclass{article}
\usepackage{/globals/de}
\usepackage{/globals/laws}
\usepackage{/globals/anki}
\ankiprefix{bundesregierung}
\ankitopic{Bundesregierung}
\ankishow
\begin{document}
\section{Die Bundesregierung} 
Die Bundesregierung ist das Staatsorgan mit der operativen Staatsleitung.

Sie nutzt \arts{62}{69} \GG und die \gesetz{GO BReg} als Rechtsquellen.
\ankinote{rechtsq}{Rechtsquellen}{Art. 62-69 GG, GO BReg}

\subsection{Funktionen}
Die Bundesregierung ist das Staatsorgan mit den meisten Gesetzesinitiativen, sie erstellt den Haushaltsplan und leitet die Exekutive der Bundesgesetzen.
\ankinote{funktionen}{Funktionen}{Gesetzesinitiativen, erstellt Haushaltsplan, Exekutive}

\subsection{Regierungsbildung}
Die Bundesregierung wird gebildet indem der Bundestag auf Vorschlag des Bundespräsidenten den Bundeskanzler wählt, der Bundesminister vorschlägt die dann vom Präsidenten ernannt werden.
\ankinote{bildung}{Bildung}{BT wählt BuKA (Vorschlag Präsident), Präsident ernennt Minister (Vorschlag BuKa)}
\end{document}